\documentclass[12pt]{article}

% ================================================
% CÀI ĐẶT TRANG & PHÔNG CHỮ TIẾNG VIỆT
% ================================================
\usepackage[margin=2cm]{geometry}
\usepackage[utf8]{inputenc}
\usepackage[T5]{fontenc}
\usepackage[vietnamese]{babel}

% ================================================
% TOÁN HỌC & HÌNH VẼ (Vừa đủ cho đề thi)
% ================================================
\usepackage{amsmath, amssymb, amsthm}
\usepackage{graphicx}
\usepackage{float}
\usepackage{tikz}
\usepackage{pgfplots}
\pgfplotsset{compat=1.18}

% ================================================
% ĐỊNH DẠNG DANH SÁCH & TRÌNH BÀY
% ================================================
\usepackage{enumitem}
\usepackage{multicol}
\usepackage{fancyhdr}
\usepackage[hidelinks]{hyperref} % Nạp sau cùng để tránh xung đột

% Thiết lập khoảng cách dòng và danh sách cho đề thi thoáng, đẹp
\setlist[enumerate]{itemsep=6pt, topsep=4pt}
\renewcommand{\baselinestretch}{1.15}

% ================================================
% CẤU HÌNH HEADER & FOOTER ĐỀ THI
% ================================================
\pagestyle{fancy}
\fancyhf{}
\fancyhead[L]{\small\textit{Bài kiểm tra thường xuyên 2 — Lớp: 2026 - L01}}
\fancyhead[R]{\small\textbf{ĐỀ CHÍNH THỨC}}
\fancyfoot[C]{Trang \thepage}
\renewcommand{\headrulewidth}{0.4pt}

% ================================================
% BẮT ĐẦU VĂN BẢN
% ================================================
\begin{document}
	
	% --- TIÊU ĐỀ ĐỀ THI CHÍNH THỨC ---
	\begin{center}
		{\Large \textbf{BÀI KIỂM TRA THƯỜNG XUYÊN 2}} \\[6pt]
		{\large \textbf{Lớp: 2026 - L01}} \\[6pt]
		\textit{Thời gian làm bài: 90 phút}
	\end{center}
	\vspace{3mm}
	\hrule
	\vspace{6mm}
	
	% ================================================
	% ĐỀ CHÍNH THỨC - PHẦN I: TRẮC NGHIỆM ĐƠN (12 CÂU)
	% ================================================
	\noindent\textbf{PHẦN I. Câu trắc nghiệm nhiều phương án lựa chọn.} Thí sinh trả lời từ câu 1 đến câu 12. Mỗi câu hỏi thí sinh chỉ chọn một phương án.
	\vspace{3mm}
	
	\begin{enumerate}[label=\textbf{Câu \arabic*.}, leftmargin=*]
		
		% CÂU 1: Vector - Oxyz (Tính toán tọa độ)
		\item Trong không gian $Oxyz$, cho hai vecto $\vec{a} = (1; 2; -1)$ và $\vec{b} = (2; -1; 3)$. Tọa độ của vecto $\vec{u} = 2\vec{a} - \vec{b}$ là:
		\begin{multicols}{4}
			\begin{enumerate}[label=\Alph*.]
				\item $\vec{u} = (0; 3; 1)$.
				\item $\vec{u} = (0; 5; -5)$.
				\item $\vec{u} = (4; 3; 1)$.
				\item $\vec{u} = (0; 5; 5)$.
			\end{enumerate}
		\end{multicols}
		
		% CÂU 2: Hàm số - Tiệm cận
		\item Đường tiệm cận ngang của đồ thị hàm số $y = \dfrac{2x - 3}{x + 1}$ có phương trình là:
		\begin{multicols}{4}
			\begin{enumerate}[label=\Alph*.]
				\item $x = -1$.
				\item $x = 2$.
				\item $y = 2$.
				\item $y = -1$.
			\end{enumerate}
		\end{multicols}
		
		% CÂU 3: Thống kê - Khoảng biến thiên
		\item Bảng dưới đây thống kê cự ly nhảy xa (đơn vị: mét) của một nhóm học sinh:
		\begin{center}
			\renewcommand{\arraystretch}{1.2}
			\begin{tabular}{|c|c|c|c|c|c|}
				\hline
				\textbf{Cự ly (m)} & $[2; 3)$ & $[3; 4)$ & $[4; 5)$ & $[5; 6)$ & $[6; 7)$ \\
				\hline
				\textbf{Số học sinh} & 5 & 12 & 18 & 10 & 3 \\
				\hline
			\end{tabular}
		\end{center}
		Khoảng biến thiên của mẫu số liệu ghép nhóm trên bằng:
		\begin{multicols}{4}
			\begin{enumerate}[label=\Alph*.]
				\item $7$.
				\item $2$.
				\item $4$.
				\item $5$.
			\end{enumerate}
		\end{multicols}
		
		% CÂU 4: Hàm số - Đơn điệu (Từ đồ thị)
		\item Cho hàm số $y = f(x)$ liên tục trên $\mathbb{R}$ và có đồ thị như hình vẽ bên dưới.
		\begin{center}
			\begin{tikzpicture}[scale=0.8]
				\begin{axis}[
					axis lines = middle,
					xlabel = {$x$}, ylabel = {$y$},
					ymin = -3, ymax = 3,
					xmin = -2.5, xmax = 2.5,
					xtick = {-1, 1}, ytick = {-2, 2},
					ticklabel style = {fill=white, inner sep=1pt},
					width=8cm, height=6.5cm,
					samples=100,
					axis line style={->}
					]
					% Đồ thị y = x^3 - 3x
					\addplot [domain=-2:2, thick, blue, smooth] {x^3 - 3*x};
					
					\draw[dashed, gray] (-1,0) -- (-1,2) -- (0,2);
					\draw[dashed, gray] (1,0) -- (1,-2) -- (0,-2);
					
					\fill (-1,2) circle (1.5pt);
					\fill (1,-2) circle (1.5pt);
					
					\node[below left] at (axis cs:0,0) {$O$};
				\end{axis}
			\end{tikzpicture}
		\end{center}
		Hàm số đã cho nghịch biến trên khoảng nào dưới đây?
		\begin{multicols}{4}
			\begin{enumerate}[label=\Alph*.]
				\item $(-1; 1)$.
				\item $(-\infty; -1)$.
				\item $(1; +\infty)$.
				\item $(-2; 2)$.
			\end{enumerate}
		\end{multicols}
		
		% CÂU 5: Vector HHKG - Biểu diễn vector
		\item Cho tứ diện $ABCD$. Gọi $M$ và $N$ lần lượt là trung điểm của các cạnh $AB$ và $CD$. Mệnh đề nào dưới đây là \textbf{đúng}?
		\begin{multicols}{2}
			\begin{enumerate}[label=\Alph*.]
				\item $\overrightarrow{MN} = \dfrac{1}{2}\left( \overrightarrow{AD} + \overrightarrow{BC} \right)$.
				\item $\overrightarrow{MN} = \dfrac{1}{2}\left( \overrightarrow{AC} + \overrightarrow{CB} \right)$.
				\item $\overrightarrow{MN} = \overrightarrow{AD} + \overrightarrow{BC}$.
				\item $\overrightarrow{MN} = \dfrac{1}{2}\left( \overrightarrow{AB} + \overrightarrow{CD} \right)$.
			\end{enumerate}
		\end{multicols}
		
		% CÂU 6: Hàm số - Max/Min (Công thức đa thức)
		\item Giá trị nhỏ nhất của hàm số $y = x^4 - 2x^2 + 3$ trên đoạn $[0; 2]$ bằng:
		\begin{multicols}{4}
			\begin{enumerate}[label=\Alph*.]
				\item $3$.
				\item $0$.
				\item $11$.
				\item $2$.
			\end{enumerate}
		\end{multicols}
		
		% CÂU 7: Vector Oxyz - Hình chiếu lên trục
		\item Trong không gian $Oxyz$, hình chiếu vuông góc của điểm $M(3; -4; 5)$ trên trục $Ox$ có tọa độ là:
		\begin{multicols}{4}
			\begin{enumerate}[label=\Alph*.]
				\item $(0; -4; 5)$.
				\item $(3; 0; 0)$.
				\item $(3; 0; 5)$.
				\item $(0; 0; 0)$.
			\end{enumerate}
		\end{multicols}
		
		% CÂU 8: Thống kê - Phương sai
		\item Thống kê thời gian (phút) giải một bài toán của 10 học sinh, ta có bảng số liệu ghép nhóm sau:
		\begin{center}
			\renewcommand{\arraystretch}{1.2}
			\begin{tabular}{|c|c|c|c|}
				\hline
				\textbf{Thời gian (phút)} & $[0; 4)$ & $[4; 8)$ & $[8; 12)$ \\
				\hline
				\textbf{Số học sinh} & 2 & 6 & 2 \\
				\hline
			\end{tabular}
		\end{center}
		Phương sai của mẫu số liệu ghép nhóm trên bằng:
		\begin{multicols}{4}
			\begin{enumerate}[label=\Alph*.]
				\item $2{,}53$.
				\item $36{,}0$.
				\item $6{,}4$.
				\item $6{,}0$.
			\end{enumerate}
		\end{multicols}
		
		% CÂU 9: Hàm số - Nhận diện hàm từ BBT
		\item Bảng biến thiên dưới đây là của hàm số nào?
		\begin{center}
			\begin{tikzpicture}[xscale=1.2, yscale=0.8]
				\draw (0,0) rectangle (8,3);
				\draw (0,1) -- (8,1);
				\draw (0,2) -- (8,2);
				\draw (1,0) -- (1,3);
				
				\node at (0.5, 2.5) {$x$};
				\node at (1.5, 2.5) {$-\infty$};
				\node at (4.5, 2.5) {$1$};
				\node at (7.5, 2.5) {$+\infty$};
				
				\node at (0.5, 1.5) {$y'$};
				\node at (3, 1.5) {$-$};
				\draw (4.4, 1) -- (4.4, 2);
				\draw (4.6, 1) -- (4.6, 2);
				\node at (6, 1.5) {$-$};
				
				\node at (0.5, 0.5) {$y$};
				\node at (1.5, 0.7) {$-1$};
				\node at (4.2, 0.2) {$-\infty$};
				\node at (4.8, 0.7) {$+\infty$};
				\node at (7.5, 0.2) {$-1$};
				
				\draw[->, thick, blue] (1.8, 0.6) -- (3.8, 0.3);
				\draw[->, thick, blue] (5.2, 0.6) -- (7.2, 0.3);
				
				\draw (4.4, 0) -- (4.4, 1);
				\draw (4.6, 0) -- (4.6, 1);
			\end{tikzpicture}
		\end{center}
		\begin{multicols}{2}
			\begin{enumerate}[label=\Alph*.]
				\item $y = \dfrac{-x+2}{x-1}$.
				\item $y = \dfrac{x+2}{x-1}$.
				\item $y = \dfrac{-x-2}{x-1}$.
				\item $y = \dfrac{-x+2}{x+1}$.
			\end{enumerate}
		\end{multicols}
		
		% CÂU 10: Vector Oxyz - Tích vô hướng (Vuông góc)
		\item Trong không gian $Oxyz$, cho hai vecto $\vec{u} = (1; m; -2)$ và $\vec{v} = (2; 3; 4)$. Tìm giá trị của tham số $m$ để hai vecto $\vec{u}$ và $\vec{v}$ vuông góc với nhau.
		\begin{multicols}{4}
			\begin{enumerate}[label=\Alph*.]
				\item $m = -2$.
				\item $m = 3$.
				\item $m = -3$.
				\item $m = 2$.
			\end{enumerate}
		\end{multicols}
		
		% CÂU 11: Hàm số - Max/Min (Phân thức 2/1)
		\item Gọi $M$ và $m$ lần lượt là giá trị lớn nhất và giá trị nhỏ nhất của hàm số $y = \dfrac{x^2 - x + 1}{x}$ trên đoạn $[1; 3]$. Tính tổng $M + m$.
		\begin{multicols}{4}
			\begin{enumerate}[label=\Alph*.]
				\item $\dfrac{4}{3}$.
				\item $\dfrac{10}{3}$.
				\item $\dfrac{7}{3}$.
				\item $\dfrac{8}{3}$.
			\end{enumerate}
		\end{multicols}
		
		% CÂU 12: Vector hình học - Góc giữa hai vecto
		\item Cho hình lập phương $ABCD.A'B'C'D'$. Số đo góc tạo bởi hai vecto $\overrightarrow{AB}$ và $\overrightarrow{A'C'}$ bằng:
		\begin{multicols}{4}
			\begin{enumerate}[label=\Alph*.]
				\item $90^\circ$.
				\item $60^\circ$.
				\item $45^\circ$.
				\item $0^\circ$.
			\end{enumerate}
		\end{multicols}
		
	\end{enumerate}
	
	\vspace{6mm}
	\begin{enumerate}[label=\textbf{Câu \arabic*.}, leftmargin=*]
		
		% CÂU 1 ĐÚNG/SAI: KHẢO SÁT ĐỒ THỊ (Hàm phân thức bậc 1/bậc 1) - Tọa độ nguyên
		\item Cho hàm số $y = f(x) = \dfrac{2x - 1}{x + 1}$ có đồ thị như hình vẽ bên dưới.
		\begin{center}
			\begin{tikzpicture}[scale=0.9]
				% Giới hạn khung vẽ đồ thị để hai cung không bị kéo quá dài tràn ra ngoài
				\begin{scope}
					\clip (-3.8, -3.2) rectangle (2.8, 4.8);
					
					% Tiệm cận đứng và tiệm cận ngang dạng nét đứt màu đỏ
					\draw[red, dashed, thick] (-1,-3.5) -- (-1,5);
					\draw[red, dashed, thick] (-4,2) -- (3,2);
					
					% Nhánh bên trái tiệm cận đứng (x < -1)
					\draw[domain=-3.8:-1.2, smooth, variable=\x, blue, very thick] plot ({\x}, {(2*\x - 1)/(\x + 1)});
					
					% Nhánh bên phải tiệm cận đứng (x > -1)
					\draw[domain=-0.8:2.8, smooth, variable=\x, blue, very thick] plot ({\x}, {(2*\x - 1)/(\x + 1)});
				\end{scope}
				
				% Vẽ hệ trục tọa độ đè lên trên cho sắc nét
				\draw[->] (-3.8, 0) -- (2.8, 0) node[right] {$x$};
				\draw[->] (0, -3.2) -- (0, 4.8) node[above] {$y$};
				\node[below left] at (0,0) {$O$};
				
				% Các điểm mốc tọa độ đặc biệt (bố trí hướng thông minh chống díu chữ)
				\fill[red] (-1,0) circle (1.5pt) node[above left, red] {$-1$};
				\fill[red] (0,2) circle (1.5pt) node[above right, red] {$2$};
				\fill (0.5,0) circle (1.5pt) node[below right, yshift=-0.5mm] {$0{,}5$};
				\fill (0,-1) circle (1.5pt) node[left] {$-1$};
			\end{tikzpicture}
		\end{center}
		Xét tính đúng hay sai của các mệnh đề sau:
		\begin{enumerate}[label=\alph*)]
			\item Hàm số luôn đồng biến trên tập xác định $\mathbb{R} \setminus \{-1\}$.
			\item Đồ thị hàm số nhận giao điểm hai đường tiệm cận $I(-1; 2)$ làm tâm đối xứng.
			\item Tiếp tuyến của đồ thị hàm số tại giao điểm của đồ thị với trục tung có hệ số góc bằng $3$.
			\item Có đúng $4$ điểm nằm trên đồ thị hàm số có cả hoành độ và tung độ đều là các số nguyên.
		\end{enumerate}
		
		\vspace{4mm}
		% CÂU 2 ĐÚNG/SAI: BÀI TOÁN OXYZ (Hình hộp, Dạng tam giác)
		\item Trong không gian $Oxyz$, cho hình hộp chữ nhật $OABC.O'A'B'C'$ có ba đỉnh nằm trên ba trục tọa độ là $A(2; 0; 0)$, $C(0; 4; 0)$ và $O'(0; 0; 3)$. (Điểm $O$ trùng với gốc tọa độ).
		\begin{center}
			\begin{tikzpicture}[scale=0.8, line join=round, line cap=round]
				% Trục Oxyz
				\draw[->, thick] (0,0) -- (-2,-1.5) node[below] {$x$};
				\draw[->, thick] (0,0) -- (6,0) node[right] {$y$};
				\draw[->, thick] (0,0) -- (0,5) node[above] {$z$};
				
				% Tọa độ các đỉnh 
				\coordinate (O) at (0,0);
				\coordinate (A) at (-1.5,-1.125); 
				\coordinate (C) at (4,0);
				\coordinate (B) at (2.5,-1.125);
				\coordinate (O') at (0,3);
				\coordinate (A') at (-1.5,1.875);
				\coordinate (C') at (4,3);
				\coordinate (B') at (2.5,1.875);
				
				% Vẽ nét khuất
				\draw[dashed, thick] (O) -- (A) (O) -- (C) (O) -- (O');
				
				% Vẽ nét thấy
				\draw[thick] (A) -- (B) -- (C) (A) -- (A') (B) -- (B') (C) -- (C') (O') -- (A') -- (B') -- (C') -- cycle;
				
				% Gắn nhãn
				\node[below left] at (O) {$O$};
				\node[left] at (A) {$A$};
				\node[below] at (B) {$B$};
				\node[below right] at (C) {$C$};
				\node[above left] at (O') {$O'$};
				\node[left] at (A') {$A'$};
				\node[above right] at (B') {$B'$};
				\node[above right] at (C') {$C'$};
			\end{tikzpicture}
		\end{center}
		Xét tính đúng hay sai của các mệnh đề sau:
		\begin{enumerate}[label=\alph*)]
			\item Tọa độ của điểm $B'$ là $(2; 4; 3)$.
			\item Vecto $\overrightarrow{OB'}$ và vecto $\overrightarrow{AC'}$ vuông góc với nhau.
			\item Trung điểm của đoạn thẳng $AB'$ nằm trên mặt phẳng tọa độ $(Oyz)$.
			\item Tam giác $OA'C$ là một tam giác cân.
		\end{enumerate}
		
	\end{enumerate}
	
	\vspace{6mm}
	% ================================================
	% PHẦN III: TRẢ LỜI NGẮN (4 CÂU)
	% ================================================
	\vspace{6mm}
	% ================================================
	% PHẦN III: TRẢ LỜI NGẮN (4 CÂU)
	% ================================================
	\noindent\textbf{PHẦN III. Câu trắc nghiệm trả lời ngắn.} Thí sinh trả lời từ câu 1 đến câu 4.
	\vspace{3mm}
	
	\begin{enumerate}[label=\textbf{Câu \arabic*.}, leftmargin=*]
		
		% CÂU 1 TRẢ LỜI NGẮN: HÀM SỐ (Tính cực đại, cực tiểu)
		\item Cho hàm số $y = 2x^3 - 3x^2 - 12x + 10$. Gọi $y_{CĐ}$ và $y_{CT}$ lần lượt là giá trị cực đại và giá trị cực tiểu của hàm số đã cho. Tính giá trị của biểu thức $P = y_{CĐ} - y_{CT}$.
		
		\vspace{2mm}
		% CÂU 2 TRẢ LỜI NGẮN: THỐNG KÊ (Tính trung vị)
		\item Khảo sát tuổi thọ (đơn vị: tháng) của 50 bóng đèn LED cùng một lô sản xuất, ta thu được mẫu số liệu ghép nhóm như sau:
		\begin{center}
			\renewcommand{\arraystretch}{1.2}
			\begin{tabular}{|c|c|c|c|}
				\hline
				\textbf{Tuổi thọ (tháng)} & $[10; 20)$ & $[20; 30)$ & $[30; 40)$ \\
				\hline
				\textbf{Số bóng đèn} & 10 & 20 & 20 \\
				\hline
			\end{tabular}
		\end{center}
		Tính số trung vị (chính là tứ phân vị thứ hai $Q_2$) của mẫu số liệu ghép nhóm trên (viết kết quả dưới dạng số thập phân).
		
		\vspace{2mm}
		% CÂU 3 TRẢ LỜI NGẮN: VECTOR HHKG (Phân tích vecto)
		\item Cho hình hộp $ABCD.A'B'C'D'$. Gọi $M$ là tâm của mặt bên $BCC'B'$. Biết rằng vecto $\overrightarrow{AM}$ được phân tích theo ba vecto cơ sở $\overrightarrow{AB}, \overrightarrow{AD}, \overrightarrow{AA'}$ dưới dạng:
		\[ \overrightarrow{AM} = x\overrightarrow{AB} + y\overrightarrow{AD} + z\overrightarrow{AA'} \quad (x, y, z \in \mathbb{R}) \]
		Tính tổng $S = x + y + z$.
		
		\vspace{2mm}
		% CÂU 4 TRẢ LỜI NGẮN: HỆ TRỤC TỌA ĐỘ OXYZ (Tìm đỉnh hình bình hành)
		\item Trong không gian $Oxyz$, cho ba điểm $A(1; 2; -1)$, $B(2; 0; 3)$ và $C(-1; 4; 2)$. Tìm tọa độ đỉnh $D(x_D; y_D; z_D)$ sao cho tứ giác $ABCD$ là một hình bình hành. Tính tổng các tọa độ của điểm $D$.
		
	\end{enumerate}
	
	\vspace{6mm}
	\vspace{6mm}
	% ================================================
	% ĐỀ CHÍNH THỨC - PHẦN IV: TỰ LUẬN (2 BÀI)
	% ================================================
	\noindent\textbf{PHẦN IV. Tự luận.} Thí sinh trình bày bài giải chi tiết từ bài 1 đến bài 2.
	\vspace{3mm}
	
	\begin{enumerate}[label=\textbf{Bài \arabic*.}, leftmargin=*]
		
		% BÀI 1: KHẢO SÁT HÀM SỐ (1.5 ĐIỂM)
		\item Cho hàm số $y = -x^3 + 3x^2 - 4$ có đồ thị là đường cong $(C)$.
		\begin{enumerate}[label=\alph*)]
			\item  Khảo sát sự biến thiên (đồng biến, nghịch biến, cực trị) và xác định tọa độ tâm đối xứng của đồ thị hàm số đã cho.
			\item  Viết phương trình tiếp tuyến của đồ thị $(C)$ tại các giao điểm của đồ thị $(C)$ với trục hoành.
		\end{enumerate}
		
		\vspace{4mm}
		% BÀI 2: VECTO TRONG KHÔNG GIAN (1.5 ĐIỂM)
		\item 
		\begin{enumerate}[label=\alph*)]
			\item  Cho hình chóp $S.ABCD$ có đáy $ABCD$ là hình bình hành tâm $O$. Chứng minh đẳng thức vecto sau:
			\[ \overrightarrow{SA} + \overrightarrow{SB} + \overrightarrow{SC} + \overrightarrow{SD} = 4\overrightarrow{SO} \]
			\item  Giả sử hình chóp ở câu trên là hình chóp tứ giác đều có tất cả các cạnh (cạnh đáy và cạnh bên) đều bằng $a$ (Hình minh họa bên cạnh). Tính giá trị của biểu thức tích vô hướng sau theo $a$:
			\[ \overrightarrow{SA} \cdot \overrightarrow{BC} \]
		\end{enumerate}
		\hfill
		\begin{center}
			\begin{tikzpicture}[scale=0.9, line join=round, line cap=round]
				% Hình chóp tứ giác đều
				\coordinate (A) at (0,0);
				\coordinate (D) at (3.5,0);
				\coordinate (B) at (-1.5,-1.5);
				\coordinate (C) at (2,-1.5);
				\coordinate (O) at (1,-0.75); 
				\coordinate (S) at (1,3.5); 
				
				% Vẽ nét khuất
				\draw[dashed, thick] (A) -- (B) (A) -- (D) (A) -- (C) (B) -- (D) (S) -- (A) (S) -- (O);
				
				% Vẽ nét thấy
				\draw[thick] (B) -- (C) -- (D) -- (S) -- (B) (S) -- (C);
				
				\fill (O) circle (1.5pt);
				
				% Gắn nhãn
				\node[left] at (A) {$A$};
				\node[below left] at (B) {$B$};
				\node[below right] at (C) {$C$};
				\node[right] at (D) {$D$};
				\node[above] at (S) {$S$};
				\node[below, yshift=-1mm] at (O) {$O$};
			\end{tikzpicture}
		\end{center}
		
	\end{enumerate}
	
\end{document}

% =========================================================================
% BẢNG TỔNG HỢP ĐÁP ÁN ĐỀ CHÍNH THỨC KTTX 2 (DÀNH CHO GIÁO VIÊN)
% =========================================================================
% 
% PHẦN I. TRẮC NGHIỆM ĐƠN
% -------------------------------------------------------------------------
% | Câu  1: B  | Câu  2: C  | Câu  3: D  | Câu  4: A  | Câu  5: A  | Câu  6: D  |
% | Câu  7: B  | Câu  8: C  | Câu  9: A  | Câu 10: D  | Câu 11: B  | Câu 12: C  |
% -------------------------------------------------------------------------
%
% PHẦN II. TRẮC NGHIỆM ĐÚNG/SAI
% -------------------------------------------------------------------------
% | CÂU 1:    a) S       |    b) Đ       |    c) Đ       |    d) Đ       |
% | CÂU 2:    a) Đ       |    b) S       |    c) S       |    d) S       |
% -------------------------------------------------------------------------
%
% PHẦN III. TRẮC NGHIỆM TRẢ LỜI NGẮN
% -------------------------------------------------------------------------
% | CÂU 1:  27     | CÂU 2:  27,5     | CÂU 3:  2        | CÂU 4:  2        |
% -------------------------------------------------------------------------
% 
% PHẦN IV. TỰ LUẬN (ĐÁP SỐ CHÍNH)
% - Bài 1a: Cực tiểu A(0; -4), Cực đại B(2; 0). Tâm đối xứng I(1; -2).
% - Bài 1b: Có 2 tiếp tuyến tại giao điểm Ox là: y = 0 và y = -9x - 9.
% - Bài 2a: Áp dụng quy tắc trung điểm với tâm đáy O: SA+SC=2SO, SB+SD=2SO.
% - Bài 2b: Tích vô hướng SA.BC = SA.AD = -a^2/2.
% =========================================================================